% ============================================================
%  01_introduction.tex --- Introduction
% ============================================================

% ==============================================================
\chapter{Introduction}
\label{ch:introduction}
% ==============================================================

% --------------------------------------------------------------
\section{From a Lagrangian to a prediction}
\label{sec:intro-motivation}
% --------------------------------------------------------------

In particle physics, we are always searching for new theories, and the key to describing a system is its Lagrangian. A theory only becomes meaningful when it is corroborated by many experiments,\supervisor{[DONE] experiments (many)} and the experimentalist who wants to test that theory needs theoretical predictions: inclusive cross sections, differential distributions and other quantities that can be measured in a particle accelerator.\supervisor{[DONE] needs theoretical predictions: inclusive cross sections, differential distributions etc.} The first step towards a bridge\supervisor{[DONE] the bridge -> the first step towards a bridge} between the theory and those measurable quantities is the set of Feynman rules, the elementary vertices from which each amplitude is assembled.

This conversion is purely mechanical and algebraic. \Cref{sec:feynrules-algorithm} describes the algorithm\supervisor{[DONE] recipe->algorithm} in detail, but we can briefly summarize it as follows: take a term of the Lagrangian, add a creation operator for each field, commute the operators, replace each derivative with a momentum\supervisor{[DONE] momentum!}, and eliminate external wave functions. Precisely because these operations are always the same, it is so prone to being implemented algorithmically. Clearly, there are many ramifications and complications: each field carries Lorentz indices, spinors, color, weak isospin, and flavor that must be contracted with the correct partner; each covariant derivative expands into a term for each gauge group under which the field transforms; each field strength of a non-Abelian group splits into three parts, so a product of $n$ of them splits into $3^n$; each exchange of two fermions costs a sign, the convention of which must be established once and never violated again.

For a Lagrangian with few terms and few fields, performing these operations manually can already be quite tedious. For the Standard Model, it is a well-known exercise but indeed error-prone as the terms are numerous. Then we have theories that are both more complex and far larger,\supervisor{[DONE] what do you mean long-lived?} such as the Standard Model Effective Field Theory (SMEFT). This is the model used as the main benchmark for this thesis, which contains 32 Lagrangian blocks and 298 independent Wilson coefficients, and generates 184 distinct reference vertices with three to six external legs. Nowadays, no one derives it manually,\supervisor{[DONE] Nowadays,..} verifies it manually, or, more precisely, modifies an operator and recalculates it manually.

With the creation of new theories that have become increasingly complex, it became necessary to automate this process. Thus, FeynRules was born, a model that uses Mathematica to calculate Feynman's rules given a theory and its Lagrangian. Years passed, and many physicists wanted to stop working with mathematics, and thus the idea for this thesis was born: to develop a toolkit capable of calculating Feynman's rules in Python.

% --------------------------------------------------------------
\section{The state of the art}
\label{sec:intro-state-of-the-art}
% --------------------------------------------------------------



The reference implementation of this idea is
\textsc{FeynRules}~\cite{Christensen:2008py,Alloul:2013bka}, a
\textsc{Mathematica} package that has been a standard entry point for building models
for collider phenomenology for over a decade. Its workflow is the one
this thesis deliberately imitates, so it is worth specifying explicitly. The \textsc{FeynRules} workflow is briefly summarized below.

First, a \textsc{FeynRules} model file is created, which has a declarative structure. The user fills
\texttt{M\$ClassesDescription} with the particle content, specifying for each field
its spin, self-conjugation, indices, and quantum numbers; they fill
\texttt{M\$Parameters} with the parameters and mixing matrices; and declare
the gauge groups in \texttt{M\$GaugeGroups}, each with its own coupling,
gauge boson, structure constants, and representations. The Lagrangian
is then written as a normal \textsc{Mathematica} expression, using a small
vocabulary of symbols: \texttt{DC[...]} for a covariant derivative,
\texttt{FS[...]} for the field-strength tensor, \texttt{del[...]} for an ordinary derivative, and \texttt{Ga[...]} for a Dirac matrix. A single call to
\texttt{FeynmanRules[...]} returns the vertices, and a further call exports
the model in the UFO format, which is used by event generators such as \textsc{MadGraph5\_aMC@NLO}.

There are two main reasons why many physicists may wish to move away from \textsc{FeynRules}. The first is that it is based entirely on \textsc{Mathematica}. This presents some inconveniences, the main one being that it does not integrate naturally with the Python-based machine-learning and analysis stack that the field now relies on for much of its other work. The second issue is computational cost. When the theory becomes complicated or the Lagrangian contains a large number of terms, the required computation time can become very high. As an example, this thesis aimed to implement the SMEFT model, which in \textsc{FeynRules} requires about half an hour of wall-clock time to export its reference vertices, as recorded in
\Cref{tab:validation-performance}.

The gap this thesis aims to fill, therefore, isn't that the problem is unsolved or that a suitable toolkit doesn't already exist—not at all.
The point is that, until recently, there was no equivalent solution available in the Python ecosystem, where much of the remaining work is taking place.

% --------------------------------------------------------------
\section{Why Python, and why now}
\label{sec:intro-why-now}
% --------------------------------------------------------------

For a long time, Python was not an adequate environment for creating such a toolkit that could perform algebraic operations and, more importantly, understand the meaning of symbolic expressions such as tensor networks with typed indices, expanded over thousands of terms.

This changed, however, with the birth of Symbolica, a computer algebra system with a core written in Rust and, more importantly, a Python interface. It provides efficient manipulation of large symbolic expressions through pattern-based rewriting and a fast multivariate polynomial engine, allowing expressions containing thousands of terms to be transformed and normalized efficiently. What Symbolica does not provide by itself is the type information needed to interpret tensor indices.

Building on Symbolica, a very promising and capable library, \textsc{Spenso}~\cite{spenso}, was developed by Lucien Huber. This library adds the missing type information: it makes an
index not simply a name, but part of a slot that also carries a representation, so that the
library itself knows that a fundamental color index can contract with its
dual and not with a Lorentz index. A complementary layer, \textsc{idenso}, provides the
additional physical identities: metric contraction, the $\mathrm{SU}(N)$ color algebra, and the Dirac algebra.

Together, these three libraries provide exactly the services that a \textsc{FeynRules}-like tool needs and that Python previously lacked.
\Cref{ch:cas} describes them in detail. This thesis is, in a nutshell, an attempt to see how far these libraries can go by further developing them.

% --------------------------------------------------------------
\section{Contributions}
\label{sec:intro-contributions}
% --------------------------------------------------------------

This thesis presents \textsc{FeynPy}, a Python toolkit that derives
tree-level Feynman rules from declarative model definitions. Its
contributions are the following.

\begin{enumerate}
    \item \textbf{A declarative model language in plain Python.} Indices,
    fields, parameters, gauge representations and gauge groups are ordinary
    Python objects, and a Lagrangian is written as an ordinary Python
    expression in a notation deliberately close to the
    \textsc{FeynRules} one
    (\Cref{ch:feynpy}).

    \item \textbf{A generic gauge compiler.} Covariant derivatives, field
    strengths, gauge-fixing terms and Faddeev--Popov ghosts are expanded
    from the model's own gauge metadata.

    \item \textbf{A contraction engine with a controlled fermion-sign
    convention.} Vertices are produced by a pruned sum over permutations of
    the external legs, the direct realisation of the canonical-quantisation
    recipe.

    \item \textbf{A canonicalisation layer that makes cross-tool comparison
    possible.} Dummy relabelling, tensor symmetries, a deterministic
    generator-ordering rewrite and a restricted use of the Jacobi identity
    reduce two independently derived expressions to a common normal form
    (\Cref{sec:internals-postprocessing,sec:validation-canonical}). 

    \item \textbf{An operator calculus on compiled Lagrangians.} Operators
    can be applied to an already compiled Lagrangian.\supervisor{[DONE] Stuff $\to$ Operations } Field
    transformations implement electroweak symmetry breaking and the
    background-field split, and graded operators implement infinitesimal
    gauge variations and BRST transformations
    (\Cref{sec:feynpy-transformations,sec:feynpy-brst}).

    \item \textbf{A model-agnostic comparison framework, and three validated
    benchmarks.} The packaged Standard Model reproduces its
    \textsc{FeynRules} export at $163/163$ vertices; the unbroken
    background-field Standard Model at $67/67$; and the Green-basis SMEFT at
    $184/184$ reference lines, of which $161$ match directly and the rest
    after documented charge-conjugation conventions (\Cref{ch:validation}).
\end{enumerate}

% --------------------------------------------------------------
\section{Scope and limitations}
\label{sec:intro-scope}
% --------------------------------------------------------------

Finally, I would like to clearly specify the scope of this work so that the results in \Cref{ch:validation} are interpreted correctly.

\textsc{FeynPy} derives Feynman rules exclusively at tree level. It does not support loop corrections, counterterms, restriction files, or export to formats such as UFO, \textsc{FeynArts}, or \textsc{CalcHEP}. Integration by parts is currently limited to scalar fields.

As a result, validation is limited to the interaction vertices and does not cover the full model metadata or numerical parameter definitions.

Within these limits, however, the result is significant: for three complete, independently written models, the vertices produced by \textsc{FeynPy} agree algebraically with those produced by \textsc{FeynRules}. The precise scope and coverage of this validation are detailed in \Cref{sec:validation-coverage}.


% --------------------------------------------------------------
\section{Outline}
\label{sec:intro-outline}
% --------------------------------------------------------------

\Cref{ch:theory} aims to provide the theoretical foundation for understanding the canonical quantization, which is the functional basis of the algorithm implemented by feynpy. The chapter will address topics such as the role of symmetry, the Lie algebra mechanism that assigns fields to their indices, the structure of gauge Lagrangians, including gauge fixing, ghosts and spontaneous symmetry breaking, the logic of effective field theory and the SMEFT basis used later, and finally the canonical quantization algorithm that transforms a Lagrangian term into a vertex.

\Cref{ch:feynpy} serves as a user's manual. It builds a model from indices and fields to a complete QCD Lagrangian with gauge fixing, extracts its vertices. Additionally it describes the additional operations available on a compiled Lagrangian: field transformations
in a physical basis, gauge variations, BRST transformations,
and the included Standard Model and SMEFT implementations.

\Cref{ch:cas} introduces \textsc{Symbolica}, \textsc{Spenso}, and \textsc{idenso} and explains the small interface object that connects them to the toolkit.

\Cref{ch:internals} explains how the code works internally. It begins by explaining how to move from declarative classes to Wick's contraction engine and canonization, and then output the expression in Symbolica.

\Cref{ch:validation} provides validation between the models implemented in FeynRules and their counterparts in FeynPY. It then reports the three benchmark comparisons, examines the charge conjugation and chirality conventions that required attention, and demonstrates, through mutation tests, that the comparison is nontrivial.

\Cref{ch:conclusion} summarizes what has been accomplished, what proved challenging, and the next steps for the work.